\section{A rendszer életciklusa}

Az előző fejezetekben külön-külön kerültek bemutatásra az adatgyűjtés, a modell betanítása és az infrastruktúra egyes komponensei.
Ez a fejezet azt mutatja be, hogyan illeszkednek ezek a részek egy egységes, üzemelő rendszerré: mi a telepítés sorrendje,
hogyan indul el a pipeline, és hogyan kerül a betanított modell éles környezetbe.

\section{A betanított modell integrálása}

A betanítási folyamat végén keletkező \texttt{best.pt} súlyfájl a \texttt{judge} Docker konténerbe kerül becsomagolásra.
A konténer induláskor betölti ezt a fájlt, és az Ultralytics YOLOv8 API-ján keresztül elérhetővé teszi azt
a Vertex AI végpont HTTP interfészén. A modellcsere folyamata a következő lépésekből áll:

\begin{enumerate}
    \item Az új \texttt{best.pt} fájlt a \texttt{judge/} könyvtárba másoljuk.
    \item A \texttt{docker-judge.yml} GitHub Action workflow automatikusan lefut, újrabuildeli a \texttt{judge} image-t,
          és feltölti az Artifact Registry-be, a commit SHA-val mint egyedi taggel.
    \item A Vertex AI végpont a következő aktiváláskor az új image-t húzza le és futtatja.
\end{enumerate}

Ez a folyamat biztosítja, hogy a modell minden frissítése verziókövetett, visszagörgetható és automatizált legyen, kézi szerverhozzáférés nélkül.

\section{Telepítési sorrend}

Az első telepítés során a következő sorrendben kell elvégezni a lépéseket:

\subsection{Cloud infrastruktúra provisioning}

Az infrastruktúra felállítása Terraform segítségével történik. A \texttt{terraform apply} parancs kiadása után
automatikusan létrejönnek a Pub/Sub topicok, a Cloud Functions, a Vertex AI végpont, a Firestore kollekciók
és az Artifact Registry repository. Az infrastruktúra változtatásai a \texttt{main} ágra való merge esetén
a CI/CD pipeline automatikusan alkalmazza, így a rendszer mindig a konfigurációval összhangban van.

\subsection{A \texttt{judge} konténer telepítése}

A \texttt{judge} Docker image buildelése és Artifact Registry-be való feltöltése szintén automatizált:
a \texttt{docker-judge.yml} GitHub Action lefuttatja a buildet és a push eseményt, amelyet a Vertex AI végpont
az első hívásakor automatikusan lehív. A konténer telepítéséhez nincs szükség manuális Docker parancsokra.

\subsection{Az edge eszköz beállítása}

A Raspberry Pi-n az alábbi három systemd szolgáltatást kell egyszeri konfigurációval engedélyezni:

\begin{itemize}
    \item \textbf{MediaMTX} - a kamerastream RTSP publikálásáért felelős.
    \item \textbf{frame-extractor} - a képkockák kinyeréséért és Pub/Sub-ra való publikálásáért felelős.
    \item \textbf{cloudflared} - a Cloudflare alagút fenntartásáért felelős, amelyen keresztül a dashboard eléri az RTSP streameket.
\end{itemize}

A szolgáltatások az eszköz bekapcsolásakor automatikusan elindulnak, emberi beavatkozás nélkül.
A \texttt{frame-extractor} konfigurációja (RTSP URL-ek, GCP projekt azonosítója, topic neve) környezeti
változókon keresztül adható meg, így ugyanaz a konténer különböző nyomtatóállomásokon is újrafelhasználható.

Mivel a képkockák időbélyege a rögzítés pillanatában keletkezik, és a felhő oldalán a Dispatcher az ennél elavultabb
képkockákat eldobja, a Pi órájának pontosnak kell lennie; ehhez az eszközön futnia kell egy idő-szinkronizáló
szolgáltatásnak (például \texttt{systemd-timesyncd} vagy NTP).

\section{Az üzemelési folyamat}

Miután a rendszer telepítve és elindítva van, a következő folyamat zajlik le automatikusan,
emberi beavatkozás nélkül (lásd \ref{fig:systemarch}. ábra):

\begin{enumerate}
    \item A Raspberry Pi-n futó \texttt{frame-extractor} szolgáltatás a három RTSP streamből
          \textbf{kameránként tíz másodpercenként egy képkockát} dolgoz fel (alapértelmezetten 0,1 fps), átméretezi és JPEG-be tömöríti azokat,
          majd JSON üzenetként publikálja a \texttt{frames-in} Pub/Sub topicra.

    \item A \textbf{Dispatcher} Cloud Function feliratkozik a \texttt{frames-in} topicra.
          Minden beérkező üzenetnél a base64-kódolt képkockát a Vertex AI végpont \texttt{/predict} végpontjára
          továbbítja OAuth2 hitelesítéssel.

    \item A \textbf{Vertex AI végponton} futó \texttt{judge} konténer elvégzi az objektumdetekciót
          a YOLOv8 modell segítségével. A téves riasztások szűrésére csak azokat a hibákat tekinti megerősítettnek, amelyek
          legalább három egymást követő képkockán át láthatók ugyanazon a kamerán. Az eredményt, a detektált hibaosztályokat,
          konfidenciaértékeket és határoló téglalapokat, minden képkockára visszapublikálja a \texttt{detections-out} topicra.

    \item Az \textbf{Alert-Manager} Cloud Function feldolgozza a detekciós eredményeket. Minden beérkező inferenciát
          naplóz a Firestore \texttt{inferences} kollekciójába, a 35\%-os konfidenciaküszöböt elérő detektálásokat pedig
          az \texttt{alerts} kollekcióba is beírja. Magas prioritású hibák (\texttt{spagetti}, \texttt{not\_sticking}, \texttt{layer\_shift}, \texttt{warping})
          esetén HTML e-mail értesítést küld a kezelőnek, egyetlen, az egész rendszerre közös 300 másodperces cooldown alkalmazásával.

    \item A \textbf{dashboard} a Firestore \texttt{onSnapshot} listenerén keresztül valós időben jeleníti meg
          az újonnan érkező riasztásokat, egyidejűleg pedig a Cloudflare alagúton keresztül élőben mutatja
          a három kamera videofolyamát.
\end{enumerate}

\section{A dashboard használata}

A webes felület megnyitása után az operátor két nézetből felügyeli a teljes nyomtatási folyamatot:

\begin{itemize}
    \item \textbf{Live streams} - a három kamera élő WebRTC videofolyamát jeleníti meg egymás mellett, a nemrég megerősített
          detekciók jelöléseivel együtt. Ezen a nézeten található a rögzítés be- és kikapcsolására szolgáló vezérlő is,
          amely csak operátori Google-fiókkal bejelentkezve használható.
    \item \textbf{Inference log} - a feldolgozott képkockák naplója: minden bejegyzéshez megjeleníti a képkocka-bélyeget és
          a hozzá tartozó detekciók táblázatát (hibaosztály, konfidenciaérték, kamera azonosítója, időbélyeg), a határoló téglalapokat
          pedig pontosan a feldolgozott képkockára rajzolva.
    \item \textbf{Állapotjelzés} - magas prioritású hiba esetén a felület piros kerettel jelez (\textit{alerting}
          állapot); költségvetési küszöb átlépésekor sárga kerettel (\textit{budget-alert} állapot).
\end{itemize}

\section{Konfigurálhatóság és skálázhatóság}

A rendszer tervezésekor fontos szempont volt, hogy könnyen adaptálható legyen különböző üzemelési körülményekhez.

\subsection{Konfidenciaküszöb}

Az Alert-Managerben beállított 35\%-os konfidenciaküszöb az alkalmazás environment változóján keresztül módosítható.
Magasabb küszöb kevesebb, de megbízhatóbb riasztást eredményez; alacsonyabb küszöb érzékenyebb, de több fals pozitívot
is hozhat. Az optimális érték meghatározása a jövőbeli tesztelési terv részeként valósul majd meg (lásd Összefoglalás és jövőbeli irányok fejezet).

\subsection{Képkockaráta és felbontás}

A \texttt{frame-extractor} \texttt{CAPTURE\_FPS} és \texttt{FRAME\_WIDTH}/\texttt{FRAME\_HEIGHT} paraméterei
lehetővé teszik a valós idejű terhelés és a Pub/Sub átviteli kapacitás közötti egyensúly megtalálását.
Gyorsabb nyomtatóknál magasabb fps értékkel érdemes dolgozni; hálózati sávszélesség korlát esetén
az alacsonyabb JPEG minőség csökkenti az átviteli méretet.

\subsection{Több nyomtatóállomás}

Az architektúra horizontálisan skálázható: további Raspberry Pi egységek és kamerakészletek az egyedi
\texttt{CAMERA\_IDS} és \texttt{RTSP\_URLS} értékek megadásával csatlakoztathatók a meglévő pipeline-hoz,
anélkül, hogy a cloud komponenseket módosítani kellene. A Pub/Sub és a Cloud Functions automatikusan skálázódnak
a megnövekedett üzenetterheléssel.

\subsection{Modellcsere}

Mivel a detekciós modell a \texttt{judge} konténertől elkülönítve, a \texttt{best.pt} fájl cseréjével frissíthető,
egy jobb modell telepítése nem igényli a teljes infrastruktúra újraépítését. Elegendő az új súlyfájlt
a \texttt{judge/} könyvtárba helyezni és a GitHub Action workflow-t lefuttatni.

\section{Ismert korlátok}

A jelenlegi implementáció az alábbi korlátokat mutatja:

\begin{itemize}
    \item A képkockákat tároló GCS bucket nyilvánosan olvasható, és nincs hozzá élettartam-szabály (\textit{lifecycle rule}),
          így a feltöltött képek korlátlanul gyűlnek. Éles üzemeltetéshez automatikus törlési szabály és szűkebb hozzáférés indokolt.
    \item Nincs külön életjel-figyelés (\textit{heartbeat}) a Raspberry Pi-re: ha a \texttt{frame-extractor} leáll, a \texttt{frames-in}
          csatorna elcsendesedik, de erről a rendszer nem küld riasztást, csak a dashboard csempéi sötétednek el.
    \item A Vertex AI végpont aktív használaton kívül is fenntart egy minimális példányt, amely napi
          $\approx$37 USD folyamatos üzemelési költséggel jár. Alacsony intenzitású monitorozáshoz
          a \textit{scale-to-zero} beállítás alkalmazása csökkentheti a kiadásokat, de növeli a hideginduló
          késleltetést (\textit{cold start}).
    \item A jelenlegi konfidenciaküszöb (35\%) empirikusan lett beállítva; a valós nyomtatási körülmények
          között végzett mérések alapján ezt finomhangolni szükséges (lásd Összefoglalás és jövőbeli irányok fejezet).
\end{itemize}
